\begin{lemma}[Balanced trunk rotation]
Let $(T,e_0,\varphi)$ be balanced of height $p=h(G)$.  For every trunk
index $1\le j\le p-1$, moving the uncolored status around the closing cycle
from $f_0=e_0$ to $f_j$ and rotating the trunk from $y_0$ to $y_j$
produces another balanced globally maximum triple.  The resulting data
satisfy the exact rotation relation \noderef{TrunkRotation}.
\end{lemma}
\begin{proof}
The closing incidence path in \noderef{NormalGlobalMaximumTriple} is
simple.  Its edge-nodes are therefore distinct, and none is $e_0$ because
they are colored while $e_0$ is uniquely uncolored.  Hence the cycle
identities $f_0,\ldots,f_{p-1}$ of \noderef{NormalCycleEdge} are all
distinct, even when two of them have the same endpoint pair.

Define $\varphi_j$ by making $f_j$ uncolored, assigning to $f_i$ the old
color of $f_{i+1}$ for $0\le i<j$, and leaving every other edge identity
unchanged.  The colored cycle identities alternate between the two trunk
colors.  At each intermediate cycle vertex, the two shifted incident
colors remain distinct.  At the initial endpoint $y_0$, the color shifted
onto $e_0$ was missing there; at the other end of the shifted segment, the
removed color becomes missing.  Thus no equal-colored adjacent pair is
created.  Away from the cycle nothing changes, so $\varphi_j$ is proper
and has exactly $f_j$ uncolored.

Order the trunk vertices cyclically as
\[
 y_j,y_{j+1},\ldots,y_{p-1},y_0,y_1,\ldots,y_{j-1},
\]
using $f_j$ first and the successive cycle identities thereafter, and
append the old post-trunk tail without changing a vertex or edge identity.
The Tashkinov condition is preserved: along the rotated trunk, the two
alternating colors are missing at its first two vertices in the rotated
coloring; in the unchanged tail, every old missing-color witness remains
somewhere in the rotated trunk because the shift only moves the missing
status of the two trunk colors around that trunk.  Every other color and
witness is unchanged.

The resulting sequence $T'$ has exactly the same vertex set and order as
$T$, hence it still attains the global Tashkinov order.  Its closing edge
is the predecessor $f_{j-1\pmod p}$, and the same alternating cycle gives
normal data $N'$ of height $p$.  These facts are precisely the vertex,
edge, color-shift, tail, and closing-edge clauses of
\noderef{TrunkRotation}.

Finally $p=h(G)$ by the original balanced hypothesis.  The post-trunk
tail is unchanged position by position, so every equality
$\varphi(e_{2r-1})=\varphi(e_{2r})$ with
$p<2r<t(G)$ is the same equality for $\varphi_j$ on $T'$.  Therefore
$N'$ satisfies \noderef{BalancedGlobalMaximumTriple} and is balanced.
\end{proof}
